\subsection{Ramification index. Residue class degree}

Let K be a field, L an extension of K and A' a valuation ring of L. As has been seen in § 1, no. 4, the ring A = K ∩ A' is a valuation ring of K and

\[
\mathfrak {m} (\mathrm{A}) = \mathfrak {m} (\mathrm{A} ^ {\prime}) \cap \mathrm{K}.
\]

If \(v'\) is a valuation associated with \(A'\) , the restriction v of \(v'\) to K is a valuation on K associated with A; the order group \(\Gamma_{v}\) of v is a subgroup of the order group \(\Gamma_{v'}\) of \(v'\) .

DEFINITION 1. The index \((\Gamma_v: \Gamma_{v'})\) is called the ramification index of \(v'\) over \(v\) (or over K) and denoted by \(e(v'/v)\) (or \(e(A'/A)\) , or sometimes \(e(L/K)\) ).

This index is a natural number or \(+\infty\) . If \(v_0'\) is a valuation equivalent to \(v'\) , \(e(v'/v)\) is also called the ramification index of \(v_0'\) over \(v\) . If \(e(v'/v) = 1\) , \(v'\) is called unramified over \(v\) .

On the other hand the residue field \(\kappa(A)\) of \(v\) is identified with a subfield of the residue field \(\kappa(A')\) of \(v'\) .

DEFINITION 2. The degree \([\kappa (\mathbf{A}^{\prime}):\kappa (\mathbf{A})]\) is called the residue class degree of \(v^{\prime}\) over \(v\) (or over K) and denoted by \(f(v^{\prime} / v)\) (or \(f(\mathbf{A}^{\prime} / \mathbf{A})\) , or sometimes \(f(\mathbf{L} / \mathbf{K})\) ).

This degree is a natural number or +∞.

LEMMA 1. Let K, K', K'' be three fields such that K ⊂ K' ⊂ K'', v'' a valuation on K'' and v and v' its restriction to K and K'. Then there are the relations:

\[
e (v ^ {\prime \prime} / v) = e (v ^ {\prime \prime} / v ^ {\prime}) e (v ^ {\prime} / v), \quad f (v ^ {\prime \prime} / v) = f (v ^ {\prime \prime} / v ^ {\prime}) f (v ^ {\prime} / v).\tag{1}
\]

This is obvious.

LEMMA 2. Let K be a field, L a finite extension of K of degree n, v' a valuation on L and v its restriction to K. Then the inequality

\[
e (v ^ {\prime} / v) f (v ^ {\prime} / v) \leqslant n\tag{2}
\]

holds; in particular \(e(v' / v)\) and \(f(v' / v)\) are finite.

Let us take natural numbers \(r\) and \(s\) respectively not greater than \(e(v'/v)\) and \(f(v'/v)\) . It suffices to show that \(rs \leqslant n\) . In view of the definition of \(r\) , there exist elements \(x_i\) of \(L\) ( \(1 \leqslant i \leqslant r\) ) such that \(v'(x_i) \not\equiv v'(x_j) \pmod{\Gamma_v}\) for \(i \neq j\) . In view of the definition of \(s\) , there exist elements \(y_k\) ( \(1 \leqslant k \leqslant s\) ) of the ring \(A'\) of \(v'\) whose canonical images \(\overline{y}_k\) in \(K(A')\) are linearly independent over \(K(A)\) ; obviously \(v'(y_k) = 0\) for all \(k\) . We shall show that the \(rs\) elements \(x_i y_k\) are linearly independent over \(K\) , which will certainly establish the inequality \(rs \leqslant n\) .

Suppose then that there exists a non-trivial linear relation of the form

\[
\sum_ {i, k} a _ {i k} x _ {i} y _ {k} = 0 \quad (a _ {i k} \in \mathbf {K}).\tag{3}
\]

Let us choose the indices \(j, m\) so that

\[
v ^ {\prime} (a _ {j m} x _ {j} y _ {m}) \leqslant v ^ {\prime} (a _ {i k} x _ {i} y _ {k})
\]

for every ordered pair \((i,k)\) ; then \(a_{jm} \neq 0\) . If \(i \neq j\) , then \(v'(a_{ik}x_iy_k) = v'(a_{jm}x_jy_m)\) is impossible for this would imply

\[
v ^ {\prime} (x _ {i}) - v ^ {\prime} (x _ {j}) = v ^ {\prime} (a _ {j m}) - v ^ {\prime} (a _ {i k}) \in \Gamma_ {v},
\]

contrary to the choice of the \(x_{i}\) . Multiplying (3) by \((a_{jm}x_j)^{-1}\) , we obtain a relation of the form

\[
\sum_ {k} b _ {k} y _ {k} + z = 0, \quad \text { where } \quad b _ {k} = \frac {a _ {j k} x _ {j}}{a _ {j m} x _ {j}} \in \mathrm{A} ^ {\prime}, \qquad z \in \mathrm{A} ^ {\prime}
\]

and \(v'(b_k) \geqslant 0, v'(z) > 0\) . Whence, in \(\kappa(A')\) , a relation of the form \(\sum_{k} \bar{b}_k \bar{y}_k = 0\) . As \(b_m = 1\) , this contradicts the hypothesis made on \(y_k\) .

PROPOSITION 1. Let K be a field, L an algebraic extension of K, \(v'\) a valuation on L, \(v\) its restriction to K and A and A' the rings of \(v\) and \(v'\) . Then \(\Gamma_{v'} / \Gamma_v\) is a torsion group and \(\kappa(A')\) is an algebraic extension of \(\kappa(A)\) .

Let \((\mathbf{L}_{\alpha})\) be the family of finite sub-extensions of L; let us write \(\Gamma_{\alpha} = v'(\mathbf{L}_{\alpha}^{*})\) . The group \(\Gamma_{v'}\) is the union of the right directed family consisting of the \(\Gamma_{\alpha}\) ; as the groups \(\Gamma_{\alpha}/\Gamma_{v}\) are finite (Lemma 2), \(\Gamma_{v'}/\Gamma_{v}\) is a torsion group. The argument is similar to prove that \(\kappa(\mathbf{A}')\) is an algebraic extension of \(\kappa(\mathbf{A})\) .

COROLLARY 1. The height of \(v'\) is equal to that of \(v\) .

This follows from Proposition 1 and the following lemma:

LEMMA 3. Let G' be a totally ordered group, G a subgroup of G' and S' (resp. S) the set of isolated subgroups of G' (resp. G). The mapping \(H' \mapsto H' \cap G\) maps S' onto S. This mapping is bijective if G'/G is a torsion group.

Clearly \(H' \in S'\) implies \(H' \cap G \in S\) . Now let \(H \in S\) ; let \(H'\) denote the set of \(x' \in G'\) such that there exists \(h \in H\) satisfying \(-h \leqslant x' \leqslant h\) ; it is immediately verified that \(H'\) is an isolated subgroup of \(G'\) ; then \(H' \cap G = H\) since H is isolated; hence the mapping \(H' \mapsto H' \cap G\) is surjective. Suppose finally that \(G'/G\) is a torsion group; let \(H_1'\) and \(H_2'\) be two isolated subgroups of \(G'\) such that \(H_1' \cap G = H_2' \cap G\) ; then, for example, \(H_1' \supset H_2'\) (cf.~§4, no. 4); then \(H_1'/H_2'\) is a totally ordered group and is isomorphic to a quotient group of \(H_1'/(H_1' \cap G)\) which itself is identified with a subgroup of \(G'/G\) ; hence \(H_1'/H_2'\) is a torsion group and therefore reduces to 0.

COROLLARY 2. For \(v'\) to be improper (resp. of height 1), it is necessary and sufficient that \(v\) be improper (resp. of height 1).

COROLLARY 3. Suppose that L is a finite extension of K. For \(v'\) to be discrete, it is necessary and sufficient that \(v\) be discrete.

If \(v'\) is discrete, \(\Gamma_v\) is isomorphic to a non-zero subgroup of \(\mathbf{Z}\) (Corollary 2) and hence to \(\mathbf{Z}\) . Conversely, if \(v\) is discrete, \(\Gamma_v\) is isomorphic to \(\mathbf{Z}\) and \(\Gamma_{v'} / \Gamma_v\) is a finite group (Lemma 2); hence \(\Gamma_{v'}\) is a finitely generated commutative group of rank 1 and torsion-free; consequently it is isomorphic to \(\mathbf{Z}\) .

